\notechapter{N-20230119}{Three Ways to Read a Triangle: Ratios, Circles, and Centers}

A triangle can be treated as three points in an affine plane, as three points on a circle, or as three vectors measured from a metric center. Each description sees different phenomena. Ratios tell us when lines meet; complex multiplication tells us when four points share a circle; vector sums explain the otherwise mysterious alignment of triangle centers. The pleasure is in choosing the viewpoint in which a hard-looking theorem becomes a one-line computation.\footnote{For the classical results and their elementary proofs see Coxeter and Greitzer, \emph{Geometry Revisited}, and R. A. Johnson, \emph{Advanced Euclidean Geometry}.}

\section{Ceva and Menelaus in Coordinates}

Let \(A,B,C\) be the vertices of a nondegenerate triangle. A point in the plane has \emph{barycentric coordinates} \((\alpha:\beta:\gamma)\), meaning
\[
 P=\frac{\alpha A+\beta B+\gamma C}
         {\alpha+\beta+\gamma}
\]
when the denominator is nonzero. In the projective completion, the triple remains meaningful even when its entries sum to zero. The coordinate \(\alpha\) vanishes exactly on the line \(BC\), and similarly for the other sidelines.

Put
\[
\begin{aligned}
 D&=(0:1-t:t)\in BC,\\
 E&=(u:0:1-u)\in CA,\\
 F&=(1-v:v:0)\in AB.
\end{aligned}
\]
For interior points, \(0<t,u,v<1\); the formulas also describe extensions of the sides. The side ratios are
\[
 \frac{BD}{DC}=\frac t{1-t},\qquad
 \frac{CE}{EA}=\frac u{1-u},\qquad
 \frac{AF}{FB}=\frac v{1-v},
\]
where ratios use directed coordinates along \(B\to C\), \(C\to A\), and \(A\to B\).

When do the three cevians \(AD,BE,CF\) meet? If their intersection is \(P=(\alpha:\beta:\gamma)\), the first cevian gives
\[
 \frac{\gamma}{\beta}=\frac t{1-t},
\]
and the others give
\[
 \frac{\alpha}{\gamma}=\frac u{1-u},
 \qquad
 \frac{\beta}{\alpha}=\frac v{1-v}.
\]
Multiplying, every coordinate cancels:
\[
 \boxed{\frac{BD}{DC}\frac{CE}{EA}\frac{AF}{FB}=1.}
\]
Conversely, if this product equals one and the denominators are nonzero, choose nonzero coordinates \(\alpha,\beta,\gamma\) satisfying the first two ratios. Their product condition forces the third, giving a common projective point. For interior side points this is the usual affine concurrence of Ceva's theorem.

Now ask the \emph{dual} question: when do the three points \(D,E,F\) lie on one line? Collinearity means the determinant of their coordinate triples is zero:
\[
 \det\begin{pmatrix}
 0&1-t&t\\
 u&0&1-u\\
 1-v&v&0
 \end{pmatrix}
 =(1-t)(1-u)(1-v)+tuv.
\]
Thus
\[
 \boxed{\frac{BD}{DC}\frac{CE}{EA}\frac{AF}{FB}=-1}
\]
when the three denominators are nonzero. This is Menelaus' theorem with the same ratio convention. What changes from \(+1\) to \(-1\) is the geometric question, not an arbitrary mnemonic about signs.

\begin{figure}[H]
\centering
\begin{tikzpicture}[scale=.82,>=Stealth]
 \begin{scope}[xshift=-3.6cm]
  \coordinate (A) at (0,3.1);
  \coordinate (B) at (-2.4,0);
  \coordinate (C) at (2.4,0);
  \coordinate (P) at (.15,1.0);
  \coordinate (D) at ($(A)!1.476190476!(P)$);
  \coordinate (E) at ($(B)!1.443959243!(P)$);
  \coordinate (F) at ($(C)!1.5872!(P)$);
  \draw[thick] (A)--(B)--(C)--cycle;
  \draw[blue!75!black] (A)--(D) (B)--(E) (C)--(F);
  \fill[blue!75!black] (P) circle (1.5pt) node[above right] {$P$};
  \foreach \p/\loc in {A/above,B/left,C/right,D/below,E/right,F/left}
   {\fill (\p) circle (1.2pt) node[\loc] {$\p$};}
  \node[font=\footnotesize] at (0,-.85) {Ceva: lines meet};
 \end{scope}
 \begin{scope}[xshift=3.6cm]
  \coordinate (A) at (0,3.1);
  \coordinate (B) at (-2.4,0);
  \coordinate (C) at (2.4,0);
  \coordinate (D) at (-3.3,0);
  \coordinate (E) at (.378,2.612);
  \coordinate (F) at (-1.30,1.42);
  \draw[thick] (A)--(B)--(C)--cycle;
  \draw[teal!75!black,thick] ($(D)!-.13!(E)$)--($(D)!1.08!(E)$);
  \foreach \p/\loc in {A/above,B/left,C/right,D/below,E/right,F/left}
   {\fill (\p) circle (1.2pt) node[\loc] {$\p$};}
  \node[font=\footnotesize] at (0,-.85) {Menelaus: points line up};
 \end{scope}
\end{tikzpicture}
\caption{The same three side points can be studied by cevians (a common solution of line equations) or by a transversal (a vanishing determinant).}
\end{figure}

Barycentric coordinates have stripped the triangle of its Euclidean metric: no lengths, angles, or perpendiculars appeared. The resulting statements are projective incidence conditions. Their expression in products of ordinary side ratios, however, assumes an affine chart and directed coordinates; an arbitrary projective transformation need not preserve each individual affine ratio.

\section{Ptolemy from Triangle Inequality}

Take four points \(A,B,C,D\) in this cyclic order on a convex quadrilateral. Ptolemy's theorem claims
\[
 AC\cdot BD=AB\cdot CD+BC\cdot AD
\]
if the quadrilateral is cyclic. But why should a product of diagonals equal a \emph{sum} of products of sides?

Represent the vertices by complex numbers \(a,b,c,d\). The elementary identity
\[
 (a-b)(c-d)+(a-d)(b-c)=(a-c)(b-d)
\]
holds for any four complex numbers. Taking absolute values gives
\[
 \boxed{|a-c||b-d|
 \le |a-b||c-d|+|a-d||b-c|.}
\]
We have obtained the \emph{Ptolemy inequality} for arbitrary four points in the Euclidean plane, without drawing any extra circle.

Equality in the triangle inequality occurs when the two summands
\((a-b)(c-d)\) and \((a-d)(b-c)\) have the same argument and orientation. For distinct points in a convex configuration, this is the cyclicity condition. Indeed their ratio
\[
 \frac{(a-b)(c-d)}{(a-d)(b-c)}
\]
is a positive real number precisely when the directed inscribed angles cut out the same circle with the stated cyclic order. Thus the cyclic case is the equality case of an elementary complex inequality.

There is an even more visual version: invert about \(A\). A circle through \(A\) becomes a line, so \(B,C,D\) turn into three collinear points \(B',C',D'\). Ptolemy's equality then becomes the addition of the two shorter segments on a line to obtain the longest one.

To see this algebraically, invert with radius \(\rho\). For any \(U,V\ne A\), inversion gives
\[
 U'V'=\frac{\rho^2 UV}{AU\cdot AV}.
\]
If \(A,B,C,D\) occur on the circle in that order, the image points \(B',C',D'\) lie in order along a line, so
\[
 B'D'=B'C'+C'D'.
\]
Multiply by \(AB\cdot AC\cdot AD/\rho^2\) to recover
\[
 AC\cdot BD=AD\cdot BC+AB\cdot CD.
\]
This is a memorable reason why Ptolemy has an addition sign: it is the segment-addition axiom after a circle is straightened.

\section{The Simson Line}

Let \(P\) lie on the circumcircle of \(ABC\). Project \(P\) perpendicularly onto the full lines \(AB,BC,CA\), obtaining \(F,D,E\). Simson's theorem says that \(D,E,F\) are collinear.

\begin{figure}[H]
\centering
\begin{tikzpicture}[scale=1.02,>=Stealth]
 \coordinate (A) at (105:2.35);
 \coordinate (B) at (215:2.35);
 \coordinate (C) at (330:2.35);
 \coordinate (P) at (40:2.35);
 \coordinate (F) at ($(A)!(P)!(B)$);
 \coordinate (D) at ($(B)!(P)!(C)$);
 \coordinate (E) at ($(C)!(P)!(A)$);
 \draw[gray!65,thick] (0,0) circle (2.35);
 \draw[thick] (A)--(B)--(C)--cycle;
 \draw[densely dashed,gray] (P)--(D) (P)--(E) (P)--(F);
 \draw[teal!75!black,line width=1.1pt] ($(D)!-.3!(F)$)--($(D)!1.3!(F)$);
 \foreach \p/\loc in {A/above,B/left,C/right,P/above right,D/below,E/right,F/left}
  {\fill (\p) circle (1.45pt) node[\loc=2pt,fill=white,inner sep=.7pt] {$\p$};}
 \node[font=\footnotesize,teal!75!black,below] at (0,-2.85) {the Simson line};
\end{tikzpicture}
\caption{Project one point of the circumcircle onto all three sidelines; the perpendicular feet unexpectedly lie on a single line.}
\end{figure}

A clean proof uses the unit circle. A similarity lets us take its center as the origin and radius as one, so the complex coordinates of \(a,b,c,p\) have absolute value one. The perpendicular foot from \(p\) to the line through \(u,v\) has complex coordinate
\[
 \boxed{h_{uv}=\frac12(p+u+v-uv\overline p).}
\]
To verify this, note that the line through \(u,v\) is
\[
 z+uv\overline z=u+v.
\]
The proposed \(h_{uv}\) satisfies the equation, and \(p-h_{uv}\) is perpendicular to the line.

Subtract two foot coordinates:
\[
 h_{ab}-h_{ac}=\frac12(b-c)(1-a\overline p),\qquad
 h_{ab}-h_{bc}=\frac12(a-c)(1-b\overline p).
\]
Since \(|p|=1\), we have \(1-a\overline p=(p-a)/p\). Therefore
\[
 \frac{h_{ab}-h_{ac}}{h_{ab}-h_{bc}}
 =\frac{(b-c)(p-a)}{(a-c)(p-b)}.
\]
This is a cross ratio of four points on the unit circle, hence real. Two complex difference vectors have real quotient exactly when they point along the same undirected line. It follows that \(h_{ab},h_{ac},h_{bc}\) are collinear. The special cases in which two feet coincide follow by continuity.

This proof identifies the hidden bridge: perpendicular projection is a \emph{metric} operation, but the final collinearity is detected by a \emph{circle cross ratio}. The same reality criterion also underlies the equality case of Ptolemy's inequality.

\section{The Euler Line from Vectors}

Let \(O\) be the circumcenter and use it as the vector origin. Write the position vectors of the vertices as \(a,b,c\), with
\[
 |a|=|b|=|c|=R.
\]
Consider the vector
\[
 h=a+b+c.
\]
The line through \(A\) in direction \(h-a=b+c\) is perpendicular to \(BC\), because
\[
 (b+c)\cdot(b-c)=|b|^2-|c|^2=0.
\]
Thus the point represented by \(h\) lies on the altitude from \(A\); the same calculation for \(B\) and \(C\) proves that \(h\) is the orthocenter \(H\).

The centroid has vector
\[
 g=\frac{a+b+c}{3}=\frac h3.
\]
Therefore \(O,G,H\) lie on one line and
\[
 \boxed{\overrightarrow{OG}=\frac13\overrightarrow{OH},
 \qquad OG:GH=1:2,}
\]
whenever \(O\ne H\). In the equilateral case, all three centers coincide. Euler's line is simply the line along a vector sum taken from the circumcenter.

The same origin gives a striking proof of Euler's relation between the incenter \(I\), inradius \(r\), circumradius \(R\), and their distance \(OI=d\). Let the side lengths opposite \(A,B,C\) be \(a_0,b_0,c_0\), and write \(s=a_0+b_0+c_0\). The incenter is the weighted average
\[
 i=\frac{a_0a+b_0b+c_0c}{s}.
\]
The elementary weighted-variance identity, valid when the weights sum to one, reads
\[
 \sum_jw_j|v_j|^2-
 \left|\sum_jw_jv_j\right|^2
 =\sum_{j<k}w_jw_k|v_j-v_k|^2.
\]
Apply it to \(v_1=a,v_2=b,v_3=c\), with weights \(a_0/s,b_0/s,c_0/s\). All three squared vertex norms are \(R^2\), and the three squared pairwise distances are \(c_0^2,a_0^2,b_0^2\). Hence
\[
 \begin{aligned}
 R^2-|i|^2
 &=\frac{a_0b_0c_0^2+
          b_0c_0a_0^2+
          c_0a_0b_0^2}{s^2}\\
 &=\frac{a_0b_0c_0}{s}.
 \end{aligned}
\]
The area \(\Delta\) satisfies
\[
 \Delta=\frac{a_0b_0c_0}{4R}
        =r\frac{s}{2}.
\]
Consequently \(a_0b_0c_0/s=2Rr\), giving
\[
 \boxed{OI^2=R^2-2Rr.}
\]
In particular \(R\ge2r\), with equality only if \(O=I\), i.e. the triangle is equilateral.

An elementary triangle-center formula has emerged from the same operation that appears in statistics as a variance decomposition. It measures how far the weighted center \(I\) lies from the equal-radius vertex circle.

\section{What changes when the allowed transformations change?}

The computations distinguish three useful geometries. Barycentric coordinates reveal concurrence and collinearity without requiring a metric. Inversion and complex coordinates reveal why cyclicity supports Ptolemy and why projected feet form a Simson line. Euclidean inner products are needed to define perpendiculars and circle centers, but they make Euler's center relations almost effortless.

The distinction matters in practice. A general affine transformation preserves the centroid and collinearity, but it need not preserve circumcenters or perpendicular feet. A projective transformation preserves incidences, but ordinary ratios on a line need not remain unchanged. A Möbius transformation preserves generalized circles and directed angles, but it does not preserve ordinary distances. Recognizing which transformations preserve a theorem lets us simplify the drawing without accidentally destroying the property we are trying to prove.
